% =====================================================================
%  minggu-01.tex  --  Slide Pertemuan Minggu ke-1
%  Pengantar Kriptografi (MPFA014) -- Prodi Matematika FMIPA UBB
%
%  Topik silabus : Pengantar & urgensi kriptografi
%  Sumber materi : buku_ajar/chapters/01 (Bab 1)
%  Sub-CPMK      : 2.1, 2.2
%
%  Kompilasi dari dalam folder ppt/:
%      pdflatex -interaction=nonstopmode minggu-01.tex   (dua kali)
% =====================================================================

\documentclass[aspectratio=169,11pt]{beamer}
% =====================================================================
%  beamer-preamble.tex  --  prasarana dan gaya slide kuliah
%
%  Mata kuliah : Pengantar Kriptografi (MPFA014)
%  Program     : Program Studi Matematika, FMIPA, Universitas Bale Bandung
%
%  Berkas ini TIDAK memuat \documentclass. Setiap berkas minggu memuatnya
%  sendiri lalu menarik berkas ini:
%
%      \documentclass[aspectratio=169,11pt]{beamer}
%      % =====================================================================
%  beamer-preamble.tex  --  prasarana dan gaya slide kuliah
%
%  Mata kuliah : Pengantar Kriptografi (MPFA014)
%  Program     : Program Studi Matematika, FMIPA, Universitas Bale Bandung
%
%  Berkas ini TIDAK memuat \documentclass. Setiap berkas minggu memuatnya
%  sendiri lalu menarik berkas ini:
%
%      \documentclass[aspectratio=169,11pt]{beamer}
%      % =====================================================================
%  beamer-preamble.tex  --  prasarana dan gaya slide kuliah
%
%  Mata kuliah : Pengantar Kriptografi (MPFA014)
%  Program     : Program Studi Matematika, FMIPA, Universitas Bale Bandung
%
%  Berkas ini TIDAK memuat \documentclass. Setiap berkas minggu memuatnya
%  sendiri lalu menarik berkas ini:
%
%      \documentclass[aspectratio=169,11pt]{beamer}
%      \input{beamer-preamble}
%      \title{...} \subtitle{...} \setWeekInfo{Minggu n}{Topik}
%
%  Berkas minggu hanya berisi \setWeekInfo dan lingkungan frame.
%  Kompilasi dijalankan DARI DALAM folder ppt/ (lihat compile-ppt.bat),
%  karena gambar buku dirujuk dengan jalur relatif ../buku_ajar/.
% =====================================================================

% ---------------------------------------------------------------- bahasa
\usepackage[utf8]{inputenc}
\usepackage[T1]{fontenc}
\usepackage[indonesian]{babel}
\usepackage{lmodern}
\usepackage{microtype}

% ------------------------------------------------------------- matematika
\usepackage{amsmath,amssymb,mathtools}

% ----------------------------------------------------------------- tabel
\usepackage{booktabs}
\usepackage{array}
\usepackage{tabularx}

% --------------------------------------------------------------- grafik
\usepackage{tikz}
\usetikzlibrary{arrows.meta,positioning,shapes.geometric,calc,fit,backgrounds}

% ---------------------------------------------------------- kode program
\usepackage{listings}
\lstdefinestyle{slidecode}{
    basicstyle=\ttfamily\scriptsize,
    keywordstyle=\color{ubbblue}\bfseries,
    commentstyle=\color{ubbgreen}\itshape,
    stringstyle=\color{ubbred},
    numbers=left,
    numberstyle=\tiny\color{gray},
    numbersep=6pt,
    frame=single,
    rulecolor=\color{ubbblue!40},
    backgroundcolor=\color{ubbblue!3},
    breaklines=true,
    showstringspaces=false,
    tabsize=4,
    columns=fullflexible,
    keepspaces=true,
}
\lstset{style=slidecode}

% =====================================================================
%  WARNA INSTITUSI
% =====================================================================
\definecolor{ubbblue}{RGB}{12, 62, 120}      % warna utama
\definecolor{ubbteal}{RGB}{0, 118, 110}      % aksen kedua
\definecolor{ubbgreen}{RGB}{22, 122, 74}     % untuk blok contoh
\definecolor{ubbred}{RGB}{168, 42, 42}       % untuk blok peringatan
\definecolor{ubbgray}{RGB}{70, 76, 84}       % teks sekunder

% =====================================================================
%  IDENTITAS MATA KULIAH
% =====================================================================
\newcommand{\CourseName}{Pengantar Kriptografi}
\newcommand{\CourseCode}{MPFA014}
\newcommand{\CourseSemester}{Semester Ganjil}
\newcommand{\CourseProgram}{Program Studi Matematika}
\newcommand{\CourseFaculty}{Fakultas Matematika dan Ilmu Pengetahuan Alam}
\newcommand{\CourseInstitution}{Universitas Bale Bandung}
\newcommand{\CourseLecturer}{Cecep Suwanda, S.Si., M.Kom.}

% Label minggu, disetel tiap berkas lewat \setWeekInfo
\newcommand{\WeekLabel}{}
\newcommand{\WeekTopic}{}
\newcommand{\setWeekInfo}[2]{%
    \renewcommand{\WeekLabel}{#1}%
    \renewcommand{\WeekTopic}{#2}%
}

% =====================================================================
%  TEMA: header, footer, warna blok
% =====================================================================
\usetheme{default}
\setbeamertemplate{navigation symbols}{}
\setbeamertemplate{blocks}[rounded][shadow=false]

\setbeamercolor{palette primary}{bg=ubbblue, fg=white}
\setbeamercolor{palette secondary}{bg=ubbgray, fg=white}
\setbeamercolor{palette tertiary}{bg=ubbteal, fg=white}

\setbeamercolor{structure}{fg=ubbblue}
\setbeamercolor{normal text}{fg=black!92}
\setbeamercolor{alerted text}{fg=ubbred}

\setbeamercolor{block title}{bg=ubbblue, fg=white}
\setbeamercolor{block body}{bg=ubbblue!6, fg=black!92}
\setbeamercolor{block title alerted}{bg=ubbred, fg=white}
\setbeamercolor{block body alerted}{bg=ubbred!6, fg=black!92}
\setbeamercolor{block title example}{bg=ubbgreen, fg=white}
\setbeamercolor{block body example}{bg=ubbgreen!6, fg=black!92}

\setbeamercolor{itemize item}{fg=ubbblue}
\setbeamercolor{itemize subitem}{fg=ubbteal}
\setbeamercolor{enumerate item}{fg=ubbblue}

\setbeamerfont{headline}{size=\scriptsize}
\setbeamerfont{footline}{size=\scriptsize}
\setbeamerfont{block title}{size=\normalsize}

% Badan blok rata kiri (bukan rata kanan-kiri): teks slide lebih mudah dibaca
% dan baris terakhir yang pendek tidak lagi memicu "Underfull \hbox".
\addtobeamertemplate{block body begin}{\raggedright}{}
\addtobeamertemplate{block alerted body begin}{\raggedright}{}
\addtobeamertemplate{block example body begin}{\raggedright}{}

% Header: nama mata kuliah dan kode MK
\setbeamertemplate{headline}{%
    \begin{beamercolorbox}[wd=\paperwidth,ht=2.6ex,dp=1.2ex,leftskip=1em,rightskip=1em]{palette primary}%
        \usebeamerfont{headline}\CourseName\ \textperiodcentered\ \CourseCode%
        \hfill \CourseProgram%
    \end{beamercolorbox}%
}

% Footer: minggu dan nomor halaman
\setbeamertemplate{footline}{%
    \begin{beamercolorbox}[wd=\paperwidth,ht=2.4ex,dp=1.2ex,leftskip=1em,rightskip=1em]{palette secondary}%
        \usebeamerfont{footline}\WeekLabel\ \textperiodcentered\ \WeekTopic%
        \hfill \insertframenumber\,/\,\inserttotalframenumber%
    \end{beamercolorbox}%
}

% Halaman judul: judul, subjudul, lalu blok identitas mata kuliah
\setbeamertemplate{title page}{%
    \begin{centering}
        \vskip1.0em
        {\usebeamerfont{title}\usebeamercolor[fg]{title}\inserttitle\par}
        \vskip0.3em
        {\usebeamerfont{subtitle}\usebeamercolor[fg]{subtitle}\insertsubtitle\par}
        \vskip1.2em
        \textcolor{ubbblue}{\rule{0.72\textwidth}{0.8pt}}
        \vskip1.2em
        {\small
            \textbf{\CourseName}\ (\CourseCode)\par
            \vskip0.4em
            \CourseProgram\par
            \CourseFaculty\par
            \CourseInstitution\par
            \vskip0.6em
            \CourseSemester\ \textperiodcentered\ \WeekLabel\par
            \vskip0.8em
            Dosen Pengampu: \CourseLecturer\par
        }
        \vskip1.0em
        \textcolor{ubbblue}{\rule{0.72\textwidth}{0.4pt}}
    \end{centering}
}

% =====================================================================
%  PERINTAH BANTU
% =====================================================================
\newcommand{\R}{\mathbb{R}}
\newcommand{\N}{\mathbb{N}}
\newcommand{\Z}{\mathbb{Z}}
\newcommand{\code}[1]{\texttt{#1}}

% Rujukan gambar dari buku: \bookfig[lebar]{NN}{namaberkas.png}
% Jalur relatif terhadap folder ppt/ (dikompilasi dari sana).
\newcommand{\bookfig}[3][0.80\linewidth]{%
    \includegraphics[width=#1]{../buku_ajar/chapters/#2/figures/#3}%
}

% Keterangan sumber gambar, sepadan dengan \sumber di buku
\newcommand{\figsumber}[1]{%
    \par\vspace{1pt}{\tiny\itshape Sumber: #1\par}%
}

% Penanda blok materi yang sering dipakai di slide
\newcommand{\slidejudul}[1]{{\large\bfseries\textcolor{ubbblue}{#1}}\par\vspace{2pt}}
\newcommand{\kunci}[1]{\textcolor{ubbred}{\textbf{#1}}}

%      \title{...} \subtitle{...} \setWeekInfo{Minggu n}{Topik}
%
%  Berkas minggu hanya berisi \setWeekInfo dan lingkungan frame.
%  Kompilasi dijalankan DARI DALAM folder ppt/ (lihat compile-ppt.bat),
%  karena gambar buku dirujuk dengan jalur relatif ../buku_ajar/.
% =====================================================================

% ---------------------------------------------------------------- bahasa
\usepackage[utf8]{inputenc}
\usepackage[T1]{fontenc}
\usepackage[indonesian]{babel}
\usepackage{lmodern}
\usepackage{microtype}

% ------------------------------------------------------------- matematika
\usepackage{amsmath,amssymb,mathtools}

% ----------------------------------------------------------------- tabel
\usepackage{booktabs}
\usepackage{array}
\usepackage{tabularx}

% --------------------------------------------------------------- grafik
\usepackage{tikz}
\usetikzlibrary{arrows.meta,positioning,shapes.geometric,calc,fit,backgrounds}

% ---------------------------------------------------------- kode program
\usepackage{listings}
\lstdefinestyle{slidecode}{
    basicstyle=\ttfamily\scriptsize,
    keywordstyle=\color{ubbblue}\bfseries,
    commentstyle=\color{ubbgreen}\itshape,
    stringstyle=\color{ubbred},
    numbers=left,
    numberstyle=\tiny\color{gray},
    numbersep=6pt,
    frame=single,
    rulecolor=\color{ubbblue!40},
    backgroundcolor=\color{ubbblue!3},
    breaklines=true,
    showstringspaces=false,
    tabsize=4,
    columns=fullflexible,
    keepspaces=true,
}
\lstset{style=slidecode}

% =====================================================================
%  WARNA INSTITUSI
% =====================================================================
\definecolor{ubbblue}{RGB}{12, 62, 120}      % warna utama
\definecolor{ubbteal}{RGB}{0, 118, 110}      % aksen kedua
\definecolor{ubbgreen}{RGB}{22, 122, 74}     % untuk blok contoh
\definecolor{ubbred}{RGB}{168, 42, 42}       % untuk blok peringatan
\definecolor{ubbgray}{RGB}{70, 76, 84}       % teks sekunder

% =====================================================================
%  IDENTITAS MATA KULIAH
% =====================================================================
\newcommand{\CourseName}{Pengantar Kriptografi}
\newcommand{\CourseCode}{MPFA014}
\newcommand{\CourseSemester}{Semester Ganjil}
\newcommand{\CourseProgram}{Program Studi Matematika}
\newcommand{\CourseFaculty}{Fakultas Matematika dan Ilmu Pengetahuan Alam}
\newcommand{\CourseInstitution}{Universitas Bale Bandung}
\newcommand{\CourseLecturer}{Cecep Suwanda, S.Si., M.Kom.}

% Label minggu, disetel tiap berkas lewat \setWeekInfo
\newcommand{\WeekLabel}{}
\newcommand{\WeekTopic}{}
\newcommand{\setWeekInfo}[2]{%
    \renewcommand{\WeekLabel}{#1}%
    \renewcommand{\WeekTopic}{#2}%
}

% =====================================================================
%  TEMA: header, footer, warna blok
% =====================================================================
\usetheme{default}
\setbeamertemplate{navigation symbols}{}
\setbeamertemplate{blocks}[rounded][shadow=false]

\setbeamercolor{palette primary}{bg=ubbblue, fg=white}
\setbeamercolor{palette secondary}{bg=ubbgray, fg=white}
\setbeamercolor{palette tertiary}{bg=ubbteal, fg=white}

\setbeamercolor{structure}{fg=ubbblue}
\setbeamercolor{normal text}{fg=black!92}
\setbeamercolor{alerted text}{fg=ubbred}

\setbeamercolor{block title}{bg=ubbblue, fg=white}
\setbeamercolor{block body}{bg=ubbblue!6, fg=black!92}
\setbeamercolor{block title alerted}{bg=ubbred, fg=white}
\setbeamercolor{block body alerted}{bg=ubbred!6, fg=black!92}
\setbeamercolor{block title example}{bg=ubbgreen, fg=white}
\setbeamercolor{block body example}{bg=ubbgreen!6, fg=black!92}

\setbeamercolor{itemize item}{fg=ubbblue}
\setbeamercolor{itemize subitem}{fg=ubbteal}
\setbeamercolor{enumerate item}{fg=ubbblue}

\setbeamerfont{headline}{size=\scriptsize}
\setbeamerfont{footline}{size=\scriptsize}
\setbeamerfont{block title}{size=\normalsize}

% Badan blok rata kiri (bukan rata kanan-kiri): teks slide lebih mudah dibaca
% dan baris terakhir yang pendek tidak lagi memicu "Underfull \hbox".
\addtobeamertemplate{block body begin}{\raggedright}{}
\addtobeamertemplate{block alerted body begin}{\raggedright}{}
\addtobeamertemplate{block example body begin}{\raggedright}{}

% Header: nama mata kuliah dan kode MK
\setbeamertemplate{headline}{%
    \begin{beamercolorbox}[wd=\paperwidth,ht=2.6ex,dp=1.2ex,leftskip=1em,rightskip=1em]{palette primary}%
        \usebeamerfont{headline}\CourseName\ \textperiodcentered\ \CourseCode%
        \hfill \CourseProgram%
    \end{beamercolorbox}%
}

% Footer: minggu dan nomor halaman
\setbeamertemplate{footline}{%
    \begin{beamercolorbox}[wd=\paperwidth,ht=2.4ex,dp=1.2ex,leftskip=1em,rightskip=1em]{palette secondary}%
        \usebeamerfont{footline}\WeekLabel\ \textperiodcentered\ \WeekTopic%
        \hfill \insertframenumber\,/\,\inserttotalframenumber%
    \end{beamercolorbox}%
}

% Halaman judul: judul, subjudul, lalu blok identitas mata kuliah
\setbeamertemplate{title page}{%
    \begin{centering}
        \vskip1.0em
        {\usebeamerfont{title}\usebeamercolor[fg]{title}\inserttitle\par}
        \vskip0.3em
        {\usebeamerfont{subtitle}\usebeamercolor[fg]{subtitle}\insertsubtitle\par}
        \vskip1.2em
        \textcolor{ubbblue}{\rule{0.72\textwidth}{0.8pt}}
        \vskip1.2em
        {\small
            \textbf{\CourseName}\ (\CourseCode)\par
            \vskip0.4em
            \CourseProgram\par
            \CourseFaculty\par
            \CourseInstitution\par
            \vskip0.6em
            \CourseSemester\ \textperiodcentered\ \WeekLabel\par
            \vskip0.8em
            Dosen Pengampu: \CourseLecturer\par
        }
        \vskip1.0em
        \textcolor{ubbblue}{\rule{0.72\textwidth}{0.4pt}}
    \end{centering}
}

% =====================================================================
%  PERINTAH BANTU
% =====================================================================
\newcommand{\R}{\mathbb{R}}
\newcommand{\N}{\mathbb{N}}
\newcommand{\Z}{\mathbb{Z}}
\newcommand{\code}[1]{\texttt{#1}}

% Rujukan gambar dari buku: \bookfig[lebar]{NN}{namaberkas.png}
% Jalur relatif terhadap folder ppt/ (dikompilasi dari sana).
\newcommand{\bookfig}[3][0.80\linewidth]{%
    \includegraphics[width=#1]{../buku_ajar/chapters/#2/figures/#3}%
}

% Keterangan sumber gambar, sepadan dengan \sumber di buku
\newcommand{\figsumber}[1]{%
    \par\vspace{1pt}{\tiny\itshape Sumber: #1\par}%
}

% Penanda blok materi yang sering dipakai di slide
\newcommand{\slidejudul}[1]{{\large\bfseries\textcolor{ubbblue}{#1}}\par\vspace{2pt}}
\newcommand{\kunci}[1]{\textcolor{ubbred}{\textbf{#1}}}

%      \title{...} \subtitle{...} \setWeekInfo{Minggu n}{Topik}
%
%  Berkas minggu hanya berisi \setWeekInfo dan lingkungan frame.
%  Kompilasi dijalankan DARI DALAM folder ppt/ (lihat compile-ppt.bat),
%  karena gambar buku dirujuk dengan jalur relatif ../buku_ajar/.
% =====================================================================

% ---------------------------------------------------------------- bahasa
\usepackage[utf8]{inputenc}
\usepackage[T1]{fontenc}
\usepackage[indonesian]{babel}
\usepackage{lmodern}
\usepackage{microtype}

% ------------------------------------------------------------- matematika
\usepackage{amsmath,amssymb,mathtools}

% ----------------------------------------------------------------- tabel
\usepackage{booktabs}
\usepackage{array}
\usepackage{tabularx}

% --------------------------------------------------------------- grafik
\usepackage{tikz}
\usetikzlibrary{arrows.meta,positioning,shapes.geometric,calc,fit,backgrounds}

% ---------------------------------------------------------- kode program
\usepackage{listings}
\lstdefinestyle{slidecode}{
    basicstyle=\ttfamily\scriptsize,
    keywordstyle=\color{ubbblue}\bfseries,
    commentstyle=\color{ubbgreen}\itshape,
    stringstyle=\color{ubbred},
    numbers=left,
    numberstyle=\tiny\color{gray},
    numbersep=6pt,
    frame=single,
    rulecolor=\color{ubbblue!40},
    backgroundcolor=\color{ubbblue!3},
    breaklines=true,
    showstringspaces=false,
    tabsize=4,
    columns=fullflexible,
    keepspaces=true,
}
\lstset{style=slidecode}

% =====================================================================
%  WARNA INSTITUSI
% =====================================================================
\definecolor{ubbblue}{RGB}{12, 62, 120}      % warna utama
\definecolor{ubbteal}{RGB}{0, 118, 110}      % aksen kedua
\definecolor{ubbgreen}{RGB}{22, 122, 74}     % untuk blok contoh
\definecolor{ubbred}{RGB}{168, 42, 42}       % untuk blok peringatan
\definecolor{ubbgray}{RGB}{70, 76, 84}       % teks sekunder

% =====================================================================
%  IDENTITAS MATA KULIAH
% =====================================================================
\newcommand{\CourseName}{Pengantar Kriptografi}
\newcommand{\CourseCode}{MPFA014}
\newcommand{\CourseSemester}{Semester Ganjil}
\newcommand{\CourseProgram}{Program Studi Matematika}
\newcommand{\CourseFaculty}{Fakultas Matematika dan Ilmu Pengetahuan Alam}
\newcommand{\CourseInstitution}{Universitas Bale Bandung}
\newcommand{\CourseLecturer}{Cecep Suwanda, S.Si., M.Kom.}

% Label minggu, disetel tiap berkas lewat \setWeekInfo
\newcommand{\WeekLabel}{}
\newcommand{\WeekTopic}{}
\newcommand{\setWeekInfo}[2]{%
    \renewcommand{\WeekLabel}{#1}%
    \renewcommand{\WeekTopic}{#2}%
}

% =====================================================================
%  TEMA: header, footer, warna blok
% =====================================================================
\usetheme{default}
\setbeamertemplate{navigation symbols}{}
\setbeamertemplate{blocks}[rounded][shadow=false]

\setbeamercolor{palette primary}{bg=ubbblue, fg=white}
\setbeamercolor{palette secondary}{bg=ubbgray, fg=white}
\setbeamercolor{palette tertiary}{bg=ubbteal, fg=white}

\setbeamercolor{structure}{fg=ubbblue}
\setbeamercolor{normal text}{fg=black!92}
\setbeamercolor{alerted text}{fg=ubbred}

\setbeamercolor{block title}{bg=ubbblue, fg=white}
\setbeamercolor{block body}{bg=ubbblue!6, fg=black!92}
\setbeamercolor{block title alerted}{bg=ubbred, fg=white}
\setbeamercolor{block body alerted}{bg=ubbred!6, fg=black!92}
\setbeamercolor{block title example}{bg=ubbgreen, fg=white}
\setbeamercolor{block body example}{bg=ubbgreen!6, fg=black!92}

\setbeamercolor{itemize item}{fg=ubbblue}
\setbeamercolor{itemize subitem}{fg=ubbteal}
\setbeamercolor{enumerate item}{fg=ubbblue}

\setbeamerfont{headline}{size=\scriptsize}
\setbeamerfont{footline}{size=\scriptsize}
\setbeamerfont{block title}{size=\normalsize}

% Badan blok rata kiri (bukan rata kanan-kiri): teks slide lebih mudah dibaca
% dan baris terakhir yang pendek tidak lagi memicu "Underfull \hbox".
\addtobeamertemplate{block body begin}{\raggedright}{}
\addtobeamertemplate{block alerted body begin}{\raggedright}{}
\addtobeamertemplate{block example body begin}{\raggedright}{}

% Header: nama mata kuliah dan kode MK
\setbeamertemplate{headline}{%
    \begin{beamercolorbox}[wd=\paperwidth,ht=2.6ex,dp=1.2ex,leftskip=1em,rightskip=1em]{palette primary}%
        \usebeamerfont{headline}\CourseName\ \textperiodcentered\ \CourseCode%
        \hfill \CourseProgram%
    \end{beamercolorbox}%
}

% Footer: minggu dan nomor halaman
\setbeamertemplate{footline}{%
    \begin{beamercolorbox}[wd=\paperwidth,ht=2.4ex,dp=1.2ex,leftskip=1em,rightskip=1em]{palette secondary}%
        \usebeamerfont{footline}\WeekLabel\ \textperiodcentered\ \WeekTopic%
        \hfill \insertframenumber\,/\,\inserttotalframenumber%
    \end{beamercolorbox}%
}

% Halaman judul: judul, subjudul, lalu blok identitas mata kuliah
\setbeamertemplate{title page}{%
    \begin{centering}
        \vskip1.0em
        {\usebeamerfont{title}\usebeamercolor[fg]{title}\inserttitle\par}
        \vskip0.3em
        {\usebeamerfont{subtitle}\usebeamercolor[fg]{subtitle}\insertsubtitle\par}
        \vskip1.2em
        \textcolor{ubbblue}{\rule{0.72\textwidth}{0.8pt}}
        \vskip1.2em
        {\small
            \textbf{\CourseName}\ (\CourseCode)\par
            \vskip0.4em
            \CourseProgram\par
            \CourseFaculty\par
            \CourseInstitution\par
            \vskip0.6em
            \CourseSemester\ \textperiodcentered\ \WeekLabel\par
            \vskip0.8em
            Dosen Pengampu: \CourseLecturer\par
        }
        \vskip1.0em
        \textcolor{ubbblue}{\rule{0.72\textwidth}{0.4pt}}
    \end{centering}
}

% =====================================================================
%  PERINTAH BANTU
% =====================================================================
\newcommand{\R}{\mathbb{R}}
\newcommand{\N}{\mathbb{N}}
\newcommand{\Z}{\mathbb{Z}}
\newcommand{\code}[1]{\texttt{#1}}

% Rujukan gambar dari buku: \bookfig[lebar]{NN}{namaberkas.png}
% Jalur relatif terhadap folder ppt/ (dikompilasi dari sana).
\newcommand{\bookfig}[3][0.80\linewidth]{%
    \includegraphics[width=#1]{../buku_ajar/chapters/#2/figures/#3}%
}

% Keterangan sumber gambar, sepadan dengan \sumber di buku
\newcommand{\figsumber}[1]{%
    \par\vspace{1pt}{\tiny\itshape Sumber: #1\par}%
}

% Penanda blok materi yang sering dipakai di slide
\newcommand{\slidejudul}[1]{{\large\bfseries\textcolor{ubbblue}{#1}}\par\vspace{2pt}}
\newcommand{\kunci}[1]{\textcolor{ubbred}{\textbf{#1}}}


\title{Pengantar dan Urgensi Kriptografi}
\subtitle{Konsep dasar, layanan keamanan, dan mengapa kriptografi penting}
\setWeekInfo{Minggu 1}{Pengantar \& urgensi kriptografi}

\begin{document}

% ----------------------------------------------------------- halaman judul
\begin{frame}[plain]
    \titlepage
\end{frame}

% ------------------------------------------------------------------ peta
\begin{frame}{Peta Materi Pertemuan Ini}
    \begin{columns}[T,onlytextwidth]
        \begin{column}{0.52\textwidth}
            \slidejudul{Enam bagian}
            \begin{enumerate}
                \item Konsep dasar komunikasi
                \item Empat masalah keamanan
                \item Apa itu kriptografi
                \item Terminologi dasar
                \item Kriptanalisis dan kriptologi
                \item Urgensi dalam teknologi informasi
            \end{enumerate}
        \end{column}
        \begin{column}{0.44\textwidth}
            \begin{block}{Referensi}
                \footnotesize
                Bab 1 buku ajar; Munir, \textit{IF4020 Kriptografi},
                ITB, 2026 (modul 01).
            \end{block}
            \begin{exampleblock}{Bobot penilaian}
                \footnotesize
                Tugas 15\%, Kuis 10\%, Praktikum 20\%,\\
                Proyek 15\%, UTS 20\%, UAS 20\%.
            \end{exampleblock}
        \end{column}
    \end{columns}
\end{frame}

% ---------------------------------------------------- capaian pembelajaran
\begin{frame}{Capaian Pembelajaran}
    \begin{block}{Sub-CPMK yang dibebankan}
        \begin{itemize}
            \item \textbf{Sub-CPMK 2.1} --- Menjelaskan definisi, terminologi,
                  sejarah, dan empat layanan kriptografi.
            \item \textbf{Sub-CPMK 2.2} --- Menjelaskan urgensi kriptografi
                  dalam teknologi informasi berdasarkan kasus kebocoran data,
                  kesalahan umum, dan tantangan masa depan.
        \end{itemize}
    \end{block}
    \vspace{0.3cm}
    \begin{exampleblock}{Setelah pertemuan ini, mahasiswa mampu}
        \begin{itemize}
            \item membedakan empat layanan keamanan dan memberi contoh pelanggarannya;
            \item menyatakan sistem kriptografi sebagai \textit{quintuple};
            \item menjelaskan Prinsip Kerckhoffs;
            \item menyebutkan kasus nyata yang menunjukkan urgensi kriptografi.
        \end{itemize}
    \end{exampleblock}
\end{frame}

% =====================================================================
\section{Konsep Dasar Komunikasi}
% =====================================================================

\begin{frame}{Pihak yang Berkomunikasi dan Salurannya}
    \begin{columns}[T,onlytextwidth]
        \begin{column}{0.56\textwidth}
            Model komunikasi kriptografi memakai tiga tokoh:
            \begin{description}
                \item[Alice] pihak yang mengirim pesan.
                \item[Bob] pihak yang menerima pesan.
                \item[Eve] penyusup (\textit{eavesdropper}) yang mengancam.
            \end{description}
            \vspace{0.2cm}
            Salurannya bersifat \kunci{publik}: siapa pun dapat mendengar,
            menyalin, atau menyisipkan pesan.
        \end{column}
        \begin{column}{0.42\textwidth}
            \centering
            \bookfig[0.95\linewidth]{01}{alur-enkripsi-dekripsi.png}
            \figsumber{Munir, \textit{IF4020 Kriptografi}, ITB.}
        \end{column}
    \end{columns}
    \vspace{0.05cm}
    \begin{alertblock}{Poin kunci}
        Kriptografi tidak menyembunyikan \emph{keberadaan} pesan, melainkan
        menyembunyikan \emph{maknanya}. Penyadap boleh melihat seluruh bit yang
        melintas, tetapi bit itu harus tidak bermakna baginya.
    \end{alertblock}
\end{frame}

\begin{frame}{Penyusup dan Kemampuannya}
    \footnotesize
    \begin{tabularx}{\linewidth}{@{}>{\bfseries}p{2.9cm} X@{}}
        \toprule
        Kemampuan & Bentuk tindakan \\
        \midrule
        Menyadap & membaca pesan yang melintas di saluran publik. \\
        \addlinespace
        Menyisipkan & memasukkan pesan sendiri ke dalam saluran. \\
        \addlinespace
        Menyamar & berpura-pura menjadi pihak lain yang sah. \\
        \addlinespace
        Membajak sesi & mengambil alih sesi komunikasi yang sudah berjalan. \\
        \addlinespace
        Menolak layanan & menghambat atau melumpuhkan layanan yang tersedia. \\
        \bottomrule
    \end{tabularx}
    \vspace{0.25cm}
    \begin{block}{Akibatnya}
        Keamanan tidak boleh bersandar pada asumsi bahwa saluran komunikasi tidak
        menarik perhatian siapa pun.
    \end{block}
\end{frame}

% =====================================================================
\section{Empat Masalah Keamanan}
% =====================================================================

\begin{frame}{Empat Pertanyaan yang Harus Dijawab}
    Ketika memakai saluran publik, muncul empat pertanyaan. Perhatikan bahwa
    keempatnya adalah pertanyaan praktis, bukan pertanyaan ahli matematika.
    \vspace{0.3cm}
    \begin{enumerate}
        \item Bagaimana Alice memastikan pesannya tidak dapat dibaca penyadap?
              \hfill\textcolor{ubbblue}{\textbf{kerahasiaan}}
        \item Bagaimana Bob memastikan pesan itu benar-benar dari Alice?
              \hfill\textcolor{ubbblue}{\textbf{otentikasi}}
        \item Bagaimana Bob memastikan pesan itu belum diubah di perjalanan?
              \hfill\textcolor{ubbblue}{\textbf{integritas}}
        \item Bagaimana Alice dicegah menyangkal telah mengirim pesan itu?
              \hfill\textcolor{ubbblue}{\textbf{nir-penyangkalan}}
    \end{enumerate}
    \vspace{0.3cm}
    \begin{exampleblock}{Catatan}
        Kriptografi tumbuh dari kebutuhan praktis, dan barulah kemudian
        dirumuskan secara matematis.
    \end{exampleblock}
\end{frame}

\begin{frame}{Kerahasiaan (\textit{Confidentiality})}
    \begin{columns}[T,onlytextwidth]
        \begin{column}{0.55\textwidth}
            Masalah ini tidak menuntut pesannya \emph{tersembunyi dari pandangan},
            melainkan \emph{tersembunyi maknanya}.
            \vspace{0.2cm}
            \begin{exampleblock}{Ilustrasi}
                Pesan ``serang fajar besok'' berubah menjadi rentetan huruf tanpa
                makna. Panjangnya bahkan boleh sama, tetapi kandungan informasinya
                tertutup.
            \end{exampleblock}
        \end{column}
        \begin{column}{0.42\textwidth}
            \centering
            \bookfig[1.0\linewidth]{01}{penyadapan-confidentiality.png}
            \figsumber{Munir, \textit{IF4020 Kriptografi}, ITB.}
        \end{column}
    \end{columns}
    \vspace{0.25cm}
    \begin{alertblock}{Prinsip Kerckhoffs}
        Kerahasiaan tidak boleh bergantung pada kerahasiaan algoritma. Musuh
        mengenal sepenuhnya sistem yang dipakai, dan \kunci{satu-satunya rahasia
        adalah kuncinya}.
    \end{alertblock}
\end{frame}

\begin{frame}{Kerckhoffs: Mengapa Algoritma Harus Publik}
    \begin{columns}[T,onlytextwidth]
        \begin{column}{0.48\textwidth}
            \begin{alertblock}{Bila algoritma dirahasiakan}
                \begin{itemize}
                    \item tidak dapat diperiksa pihak luar;
                    \item tidak dapat distandarkan;
                    \item runtuh seketika begitu satu orang membocorkannya.
                \end{itemize}
            \end{alertblock}
        \end{column}
        \begin{column}{0.48\textwidth}
            \begin{exampleblock}{Bila keamanan bergantung pada kunci}
                \begin{itemize}
                    \item membocorkan algoritma tidak merugikan siapa pun;
                    \item penggantian kunci cukup dengan membangkitkan kunci baru.
                \end{itemize}
            \end{exampleblock}
        \end{column}
    \end{columns}
    \vspace{0.35cm}
    Prinsip ini dirumuskan Auguste Kerckhoffs pada abad ke-19, lalu ditegaskan
    kembali oleh Claude Shannon dalam bentuk yang lebih tajam.
\end{frame}

\begin{frame}{Otentikasi, Integritas, dan Nir-Penyangkalan}
    \begin{block}{Otentikasi}
        Pada saluran publik tidak ada yang dapat memastikan identitas secara
        otomatis: alamat pengirim surel dapat dipalsukan, nomor telepon dapat
        dipalsukan, alamat sumber paket data dapat tidak benar. Kerahasiaan dan
        otentikasi adalah dua masalah yang berbeda.
    \end{block}
    \begin{block}{Integritas data}
        Stasiun cuaca mengirim $33^{\circ}$C, pihak tak berwenang mengubahnya
        menjadi $32^{\circ}$C. Kerahasiaan tidak dilanggar --- suhu boleh diketahui
        siapa saja --- tetapi integritasnya rusak.
    \end{block}
    \begin{block}{Nir-penyangkalan}
        Mencegah pengirim menyangkal pesan yang sudah dikirimnya. Menjadi syarat
        mutlak bagi transaksi elektronik yang mengikat secara hukum.
    \end{block}
\end{frame}

\begin{frame}{Ringkasan Empat Layanan}
    \footnotesize
    \begin{tabularx}{\linewidth}{@{}>{\bfseries}p{3.0cm} p{5.2cm} X@{}}
        \toprule
        Layanan & Pertanyaan yang dijawab & Contoh pelanggaran \\
        \midrule
        Kerahasiaan & Apakah isi pesan tertutup bagi penyadap? & Penyadapan lalu lintas tanpa enkripsi. \\
        \addlinespace
        Otentikasi & Apakah pengirim benar-benar pihak yang diklaim? & Surel dengan alamat pengirim palsu. \\
        \addlinespace
        Integritas & Apakah pesan belum diubah di perjalanan? & Pembacaan suhu $33^\circ$C diubah menjadi $32^\circ$C. \\
        \addlinespace
        Nir-penyangkalan & Dapatkah pengirim menyangkal pesannya? & Sengketa transaksi tanpa tanda tangan digital. \\
        \bottomrule
    \end{tabularx}
    \vspace{0.3cm}
    \begin{alertblock}{Penting}
        Keempatnya tidak dapat dipenuhi oleh satu mekanisme tunggal. Setiap
        layanan menuntut mekanisme tersendiri.
    \end{alertblock}
\end{frame}

% =====================================================================
\section{Apa Itu Kriptografi}
% =====================================================================

\begin{frame}{Definisi dan Etimologi}
    \begin{block}{Definisi}
        Kriptografi adalah ilmu dan seni menjaga keamanan pesan melalui teknik
        matematika.
    \end{block}
    \vspace{0.2cm}
    \begin{columns}[T,onlytextwidth]
        \begin{column}{0.46\textwidth}
            \begin{exampleblock}{Etimologi}
                \textit{crypt\'os} = rahasia\\
                \textit{gr\'aphein} = tulisan
            \end{exampleblock}
        \end{column}
        \begin{column}{0.50\textwidth}
            \begin{alertblock}{Yang dijaga}
                Bukan keberadaan pesan, melainkan maknanya --- melalui perubahan
                representasi pesan, bukan penyembunyian keberadaannya.
            \end{alertblock}
        \end{column}
    \end{columns}
\end{frame}

\begin{frame}{Layanan Kriptografi}
    Kriptografi menyediakan lebih dari sekadar kerahasiaan:
    \begin{itemize}
        \item kerahasiaan (\textit{confidentiality});
        \item integritas data (\textit{integrity});
        \item otentikasi (\textit{authentication});
        \item nir-penyangkalan (\textit{non-repudiation}).
    \end{itemize}
    \vspace{0.3cm}
    \begin{block}{Empat layanan, empat masalah}
        Setiap layanan menjawab satu masalah keamanan komunikasi pada
        bagian sebelumnya.
    \end{block}
\end{frame}

\begin{frame}{Tiga Jenis Algoritma Kriptografi}
    \begin{columns}[T,onlytextwidth]
        \begin{column}{0.32\textwidth}
            \centering
            \bookfig[0.92\linewidth]{01}{kriptografi-simetri.png}\\
            {\footnotesize\textbf{Simetri}\\ $k_{\text{enkripsi}} = k_{\text{dekripsi}}$}
        \end{column}
        \begin{column}{0.32\textwidth}
            \centering
            \bookfig[0.92\linewidth]{01}{kriptografi-nirsimetri.png}\\
            {\footnotesize\textbf{Nir-simetri}\\ $k_{\text{enkripsi}} \ne k_{\text{dekripsi}}$}
        \end{column}
        \begin{column}{0.30\textwidth}
            \vspace{0.4cm}
            \begin{block}{Fungsi hash}
                \footnotesize
                Tanpa kunci, berpanjang tetap, dan \kunci{tidak dapat dibalik}.
            \end{block}
        \end{column}
    \end{columns}
    \vspace{0.35cm}
    \begin{alertblock}{Catatan}
        Fungsi hash sengaja melanggar syarat $D_k(E_k(m)) = m$ --- dan justru
        pelanggaran itulah yang menjadi gunanya.
    \end{alertblock}
\end{frame}

\begin{frame}{Dua Aplikasi Utama Enkripsi}
    \begin{columns}[T,onlytextwidth]
        \begin{column}{0.48\textwidth}
            \begin{block}{\textit{Encryption at rest}}
                Dokumen yang tersimpan di dalam penyimpanan.\\
                \footnotesize Contoh: enkripsi basis data dan cakram keras.
            \end{block}
        \end{column}
        \begin{column}{0.48\textwidth}
            \begin{block}{\textit{Encryption in motion}}
                Pesan yang sedang dikirim melalui jaringan.\\
                \footnotesize Contoh: pesan pada aplikasi perpesanan.
            \end{block}
        \end{column}
    \end{columns}
    \vspace{0.35cm}
    Sebagian besar kebocoran data terjadi bukan karena algoritma dipecahkan,
    melainkan karena \kunci{pengabaian encryption at rest} dan pengelolaan kunci
    yang buruk.
\end{frame}

\begin{frame}{Sejarah Singkat}
    \begin{itemize}
        \item \textbf{Abad ke-9} --- \textbf{Al-Kindi} merintis kriptanalisis melalui
              \emph{analisis frekuensi}: huruf yang paling sering muncul pada
              cipherteks cenderung mewakili huruf yang paling sering muncul pada
              bahasa aslinya.
        \item Perkembangan berlanjut melalui cipher klasik (Caesar, Vigen\`ere,
              Playfair, Hill, Enigma) yang dibahas pada Bab 4.
        \item \textbf{Abad ke-20} --- Shannon meletakkan dasar teori informasi;
              lahirlah kriptografi modern yang bekerja pada bit, bukan huruf.
        \item \textbf{1976--1977} --- Diffie--Hellman dan RSA membuka era
              kriptografi kunci-publik.
    \end{itemize}
    \vspace{0.3cm}
    \begin{exampleblock}{Benang merah}
        Setiap algoritma baru melahirkan kriptanalisis baru, dan setiap
        kriptanalisis baru menuntut algoritma yang lebih kuat. Bidang ini bergerak
        karena tarik-menarik itu.
    \end{exampleblock}
\end{frame}

% =====================================================================
\section{Terminologi Dasar}
% =====================================================================

\begin{frame}{Sistem Kriptografi sebagai \textit{Quintuple}}
    Sistem kriptografi dinyatakan sebagai
    $(\mathcal{E}, \mathcal{D}, \mathcal{M}, \mathcal{K}, \mathcal{C})$ dengan:
    \vspace{0.15cm}
    \begin{columns}[T,onlytextwidth]
        \begin{column}{0.48\textwidth}
            \begin{itemize}
                \item $\mathcal{M}$ --- himpunan plainteks;
                \item $\mathcal{K}$ --- himpunan kunci;
                \item $\mathcal{C}$ --- himpunan cipherteks;
            \end{itemize}
        \end{column}
        \begin{column}{0.48\textwidth}
            \begin{itemize}
                \item $\mathcal{E}$ --- $E: \mathcal{M} \times \mathcal{K} \to \mathcal{C}$;
                \item $\mathcal{D}$ --- $D: \mathcal{C} \times \mathcal{K} \to \mathcal{M}$.
            \end{itemize}
        \end{column}
    \end{columns}
    \vspace{0.3cm}
    \begin{alertblock}{Syarat yang harus berlaku}
        \[
            D_k\bigl(E_k(m)\bigr) = m
            \qquad \text{untuk setiap } k \in \mathcal{K},\ m \in \mathcal{M}.
        \]
        Inilah yang membedakan cipher dari sekadar pengacak.
    \end{alertblock}
\end{frame}

\begin{frame}{Contoh: Caesar Cipher sebagai \textit{Quintuple}}
    \begin{itemize}
        \item $\mathcal{M} = \{$huruf-huruf alfabet$\}$
        \item $\mathcal{K} = \{k \mid k$ bilangan bulat dan $0 \le k \le 25\}$
        \item $\mathcal{E} = \{E_k \mid E_k(m) = (m + k) \bmod 26\}$
        \item $\mathcal{D} = \{D_k \mid D_k(c) = (c - k) \bmod 26\}$
        \item $\mathcal{C} = \mathcal{M}$
    \end{itemize}
    \vspace{0.3cm}
    \begin{exampleblock}{Perhatikan}
        Pada Caesar cipher, $\mathcal{C} = \mathcal{M}$. Sifat ini \emph{tidak}
        berlaku umum: pada cipher modern, cipherteks biasanya berupa deretan bit
        yang panjangnya justru dipaksa berbeda dari plainteksnya.
    \end{exampleblock}
\end{frame}

\begin{frame}{Miskonsepsi Umum}
    \footnotesize
    \begin{columns}[T,onlytextwidth]
        \begin{column}{0.52\textwidth}
            Pemahaman istilah yang longgar melahirkan kekeliruan yang berulang:
            \begin{itemize}
                \item menganggap enkripsi sama dengan \textit{encoding};
                \item menganggap enkripsi sama dengan \textit{hashing};
                \item menganggap enkripsi sama dengan tokenisasi;
                \item tidak memahami perbedaan simetri dan nir-simetri.
            \end{itemize}
        \end{column}
        \begin{column}{0.46\textwidth}
            \centering
            \bookfig[1.0\linewidth]{01}{encoding-hashing-encryption.png}
            \figsumber{Munir, \textit{IF4020 Kriptografi}, ITB.}
        \end{column}
    \end{columns}
    \begin{alertblock}{Akibatnya}
        Menyamakan kelima hal itu --- \textit{encoding}, \textit{hashing}, enkripsi,
        tokenisasi, dan tanda tangan digital --- menimbulkan \kunci{rasa aman yang
        semu}.
    \end{alertblock}
\end{frame}

% =====================================================================
\section{Kriptanalisis dan Kriptologi}
% =====================================================================

\begin{frame}{Kriptografi, Kriptanalisis, dan Kriptologi}
    \begin{columns}[T,onlytextwidth]
        \begin{column}{0.50\textwidth}
            \begin{block}{Kriptografi}
                Merancang algoritma agar pesan aman.
            \end{block}
            \begin{alertblock}{Kriptanalisis}
                Memecahkan cipherteks tanpa memiliki kunci. Dirintis Al-Kindi
                pada abad ke-9.
            \end{alertblock}
            \begin{exampleblock}{Kriptologi}
                Payung yang mencakup kriptografi \emph{dan} kriptanalisis.
            \end{exampleblock}
        \end{column}
        \begin{column}{0.46\textwidth}
            \centering
            \bookfig[1.0\linewidth]{01}{kriptologi-cabang.png}
            \figsumber{Munir, \textit{IF4020 Kriptografi}, ITB.}
        \end{column}
    \end{columns}
\end{frame}

% =====================================================================
\section{Urgensi dalam Teknologi Informasi}
% =====================================================================

\begin{frame}{Kebocoran Data di Indonesia (1/2)}
    \footnotesize
    \begin{tabularx}{\linewidth}{@{}>{\raggedright\arraybackslash}p{3.0cm} >{\raggedright\arraybackslash}p{2.4cm} >{\raggedright\arraybackslash}X@{}}
        \toprule
        Peristiwa & Skala & Keterangan \\
        \midrule
        BPJS Kesehatan (Mei 2021) & 279 juta & NIK, nama, alamat, nomor telepon, surel, foto; dijual di Raid Forums dengan harga 0,15 BTC, saat itu setara Rp70--80 juta. \\
        \addlinespace
        Tokopedia (2020) & 91 juta pengguna + 7 juta \textit{merchant} & basis data pengguna. \\
        \addlinespace
        KPU (2020) & 2,3 juta & data pemilih Pemilu 2014. \\
        \addlinespace
        Pasien Covid-19 (2020) & 250 ribu & data pasien. \\
        \bottomrule
    \end{tabularx}
    \vspace{0.2cm}
    \begin{alertblock}{Yang perlu dicermati}
        Harga jualnya sangat murah dibandingkan kerugian yang ditimbulkan. Bagi
        penyerang, biaya memperoleh data semacam itu memang rendah bila sejak awal
        tidak dilindungi.
    \end{alertblock}
\end{frame}

\begin{frame}{Kebocoran Data di Indonesia (2/2)}
    \footnotesize
    \begin{itemize}
        \item \textbf{Kementerian Pertahanan} --- 1,64 terabita data.
        \item \textbf{Bank Syariah Indonesia} (Mei 2023) --- 1,5 TB data nasabah,
              dibobol kelompok peretas LockBit 3.0.
        \item Dugaan kebocoran data pemilih Pemilu 2014 dan kebocoran sertifikat
              vaksin.
    \end{itemize}
    \vspace{0.15cm}
    \begin{alertblock}{Kesimpulan yang penting}
        Persoalannya bukan kecelakaan tunggal, melainkan \kunci{kelemahan yang
        sistemik}. Umumnya kasus itu \emph{bukan} kegagalan memecahkan algoritma
        yang kuat, melainkan kunci yang tidak dikelola dengan benar, basis data
        yang tersimpan tanpa enkripsi, atau kredensial yang bocor.
    \end{alertblock}
    \vspace{0.15cm}
    \begin{block}{Konsekuensi regulasi}
        Masyarakat berhak mengetahui tujuan penggunaan datanya dan meminta
        penghapusan; pengelola wajib memitigasi kebocoran dan dapat dikenai sanksi.
        Tanpa enkripsi yang memadai, mitigasi hanya bersifat administratif.
    \end{block}
\end{frame}

\begin{frame}{Penyadapan Komunikasi}
    \begin{itemize}
        \item \textbf{Kedutaan Besar RI di Yangon, Myanmar} --- dinilai melanggar
              Konvensi Wina karena mengganggu fasilitas kedutaan dan menyangkut
              kerahasiaan.
        \item \textbf{2013} --- dugaan penyadapan telepon Presiden Susilo Bambang
              Yudhoyono oleh pemerintah Australia; Indonesia memanggil pulang
              dubesnya dari Canberra untuk berkonsultasi.
        \item \textbf{WikiLeaks} --- 400 ribu dokumen Perang Irak (Oktober 2010) dan
              kawat diplomatik Amerika Serikat (November 2010).
    \end{itemize}
    \vspace{0.3cm}
    \begin{alertblock}{Pelajaran}
        Penyadapan tidak terbatas pada pelaku kejahatan --- negara pun
        melakukannya, dan sasarannya dapat berupa saluran komunikasi tingkat
        tertinggi. Begitu data rahasia tersimpan dalam bentuk yang dapat dibaca,
        kebocorannya hanya menunggu waktu.
    \end{alertblock}
\end{frame}

\begin{frame}{Keamanan Transaksi Elektronik}
    Kriptografi dipakai untuk melindungi transaksi elektronik, baik data yang
    disimpan maupun pesan yang dikirim.
    \vspace{0.25cm}
    \begin{columns}[T,onlytextwidth]
        \begin{column}{0.48\textwidth}
            \begin{block}{Data yang disimpan}
                Data transaksi dijaga kerahasiaannya di dalam penyimpanan.
            \end{block}
        \end{column}
        \begin{column}{0.48\textwidth}
            \begin{block}{Pesan yang dikirim}
                Pesan transaksi dijaga kerahasiaannya saat melintas jaringan.
            \end{block}
        \end{column}
    \end{columns}
    \vspace{0.3cm}
    Tanpa perlindungan itu, kepercayaan publik terhadap layanan elektronik
    sulit dibangun.
\end{frame}

\begin{frame}{Lembaga Terkait Kriptografi di Indonesia}
    \begin{enumerate}
        \item \textbf{Badan Siber dan Sandi Negara (BSSN)} --- penggabungan Lembaga
              Sandi Negara (Lemsaneg) dengan Direktorat Jenderal Aplikasi
              Informatika (Aptika) Kementerian Komunikasi dan Informatika.
        \item \textbf{Politeknik Siber dan Sandi Negara (PSSN)} ---
              menyelenggarakan pendidikan vokasi di bidang persandian.
    \end{enumerate}
    \vspace{0.3cm}
    \begin{exampleblock}{Museum Sandi, Yogyakarta}
        Beralamat di Jalan Faridan Muridan Noto Nomor 21, Kota Baru, Yogyakarta.
        Dikatakan sebagai satu-satunya museum sandi di Indonesia, bahkan di dunia.
        Koleksinya mencakup mesin sandi SN-101.
    \end{exampleblock}
\end{frame}

\begin{frame}{Tantangan Masa Depan}
    \begin{columns}[T,onlytextwidth]
        \begin{column}{0.54\textwidth}
            \begin{itemize}
                \item \textbf{Kriptografi ringan} (\textit{lightweight
                      cryptography}) --- hemat daya, memori, dan siklus komputasi,
                      layak dipasang pada perangkat IoT.
                \item \textbf{Kriptografi pasca-kuantum} --- algoritma Shor
                      mengancam RSA dan ECC, sehingga dikembangkan keluarga baru
                      berbasis \textit{hash}, kode, multivariat, dan kisi.
                \item \textbf{Kriptografi berbasis kecerdasan buatan} ---
                      mempercepat kriptanalisis sekaligus merancang primitif baru.
                \item \textbf{\textit{Selective encryption}} --- menyandikan hanya
                      bagian terpenting aliran video.
            \end{itemize}
        \end{column}
        \begin{column}{0.44\textwidth}
            \centering
            \bookfig[1.0\linewidth]{01}{lightweight-cryptography-iot.png}
            \figsumber{Munir, \textit{IF4020 Kriptografi}, ITB.}
        \end{column}
    \end{columns}
\end{frame}

% =====================================================================
\section{Rangkuman dan Latihan}
% =====================================================================

\begin{frame}{Rangkuman (1/2)}
    \footnotesize
    \begin{itemize}
        \item Kriptografi adalah ilmu dan seni menjaga keamanan pesan melalui
              teknik matematika; dari \textit{crypt\'os} dan \textit{gr\'aphein}.
        \item Dalam model komunikasi kriptografi, pengirim ditokohkan sebagai
              \textbf{Alice} dan penerima sebagai \textbf{Bob}; pihak ketiga yang
              mengancam disebut penyusup.
        \item Empat masalah keamanan komunikasi pada saluran publik: kerahasiaan,
              otentikasi, integritas data, dan nir-penyangkalan.
        \item Sistem kriptografi dimodelkan sebagai \textit{quintuple}
              $(\mathcal{E},\mathcal{D},\mathcal{M},\mathcal{K},\mathcal{C})$
              dengan syarat $D_k(E_k(m)) = m$ untuk setiap kunci yang sah.
        \item Menurut \textbf{Prinsip Kerckhoffs}, algoritma boleh dan harus
              publik; satu-satunya rahasia adalah kunci.
    \end{itemize}
\end{frame}

\begin{frame}{Rangkuman (2/2)}
    \footnotesize
    \begin{itemize}
        \item Tiga keluarga algoritma: simetri, nir-simetri, dan fungsi hash.
        \item Enkripsi dipakai untuk \textit{encryption at rest} dan
              \textit{encryption in motion}.
        \item Kriptanalisis adalah upaya memecahkan cipherteks tanpa kunci,
              dirintis Al-Kindi pada abad ke-9 melalui analisis frekuensi;
              kriptologi mencakup kriptografi sekaligus kriptanalisis.
        \item \textit{Encoding}, \textit{hashing}, enkripsi, tokenisasi, dan tanda
              tangan digital adalah lima hal yang berbeda.
        \item Urgensi nyata: kebocoran data (BPJS Kesehatan, Tokopedia, KPU,
              Kemenhan, BSI), penyadapan diplomatik, dan pengamanan transaksi
              elektronik.
        \item Di Indonesia ditangani BSSN dan PSSN; Museum Sandi di Yogyakarta
              menjadi sarana pelestarian.
        \item Tantangan ke depan: kriptografi ringan, pasca-kuantum, berbasis AI,
              dan \textit{selective encryption}.
    \end{itemize}
\end{frame}

\begin{frame}{Latihan}
    \footnotesize
    \begin{enumerate}
        \item Jelaskan perbedaan \textit{confidentiality} dan \textit{integrity}
              memakai contoh pengiriman nilai ujian mahasiswa.
        \item Mengapa Prinsip Kerckhoffs menyatakan algoritma harus publik dan
              hanya kunci yang rahasia? Apa risikonya bila algoritma juga
              dirahasiakan?
        \item Sebuah perusahaan menyimpan basis data pelanggan dalam berkas
              Base64, dan manajemennya mengklaim data itu ``sudah terenkripsi''.
              Bantah klaim itu dengan menunjukkan apa yang dapat dilakukan
              penyerang dalam hitungan detik.
        \item Mengapa empat layanan kriptografi tidak dapat dipenuhi satu
              mekanisme tunggal? Beri satu contoh ketika kerahasiaan terpenuhi
              tetapi integritas tidak.
        \item Bandingkan kriptografi simetri dan nir-simetri dari jumlah kunci
              yang dikelola dan kecepatan komputasi.
        \item Tuliskan Caesar cipher sebagai
              $(\mathcal{E},\mathcal{D},\mathcal{M},\mathcal{K},\mathcal{C})$.
        \item Sebuah tim mengamankan komunikasi 200 cabang bank. Hitung kunci yang
              diperlukan pada kriptografi simetri, dan pasangan kunci pada
              kriptografi nir-simetri.
    \end{enumerate}
\end{frame}

\begin{frame}{Checklist Pencapaian Kompetensi}
    \footnotesize
    \begin{tabularx}{\linewidth}{@{}X c@{}}
        \toprule
        Indikator Kompetensi & Tercapai \\
        \midrule
        Saya dapat menjelaskan definisi dan terminologi dasar kriptografi & $\square$ \\
        \hline
        Saya dapat membedakan 4 layanan utama (kerahasiaan, otentikasi, integritas, nir-penyangkalan) & $\square$ \\
        \hline
        Saya dapat membedakan algoritma simetri dan nir-simetri & $\square$ \\
        \hline
        Saya dapat menjelaskan Prinsip Kerckhoffs beserta konsekuensinya & $\square$ \\
        \hline
        Saya dapat memberikan contoh kasus urgensi kriptografi di dunia TI & $\square$ \\
        \bottomrule
    \end{tabularx}
    \vspace{0.35cm}
    \begin{block}{Pertemuan berikutnya}
        \textbf{Minggu 2} --- Layanan keamanan dan serangan
        (Bab 2): jenis-jenis serangan, model ancaman, dan kriptanalisis.
    \end{block}
\end{frame}

\end{document}
